\documentclass[11pt]{article}
% Version française de notes.tex, note associée à article-fr.tex ; la version anglaise est notes.tex. Compilée avec XeLaTeX :
% cd docs/bridges && xelatex notes-fr.tex && xelatex notes-fr.tex
\usepackage{polyglossia}\setmainlanguage{french}\setotherlanguage{english}
\newcommand{\thmname}{Théorème}\newcommand{\corname}{Corollaire}\newcommand{\propname}{Proposition}
\newcommand{\lemname}{Lemme}\newcommand{\bridgename}{Pont}\newcommand{\defname}{Définition}
\newcommand{\exname}{Exemple}\newcommand{\exoname}{Exercice}\newcommand{\remname}{Remarque}
% Shared preamble of docs/bridges/article.tex (English) and article-fr.tex (French).
% Built with XeLaTeX: cd docs/bridges && xelatex article.tex && xelatex article.tex
\usepackage{fontspec}
\IfFontExistsTF{Charter}{\setmainfont{Charter}}{\IfFontExistsTF{XCharter}{\setmainfont{XCharter}}{}}
\setmonofont{DejaVu Sans Mono}[Scale=0.8]
\usepackage{microtype}
\usepackage[a4paper,margin=27mm]{geometry}
\usepackage{setspace}\setstretch{1.1}
\usepackage{amsmath,amssymb,amsthm}
\usepackage{mathtools}
\usepackage[most]{tcolorbox}
\usepackage{xcolor}
\usepackage{tikz}
\usetikzlibrary{arrows.meta,positioning,calc,cd}
\usepackage{booktabs,array,tabularx}
\usepackage{fancyvrb}
\usepackage[hidelinks]{hyperref}
\usepackage{xurl}
\usepackage{enumitem}
\setlist{nosep,leftmargin=*}
\setlength{\emergencystretch}{2em}
\setlength{\parskip}{0.35em}
\numberwithin{equation}{section}

\newcommand{\Set}{\mathbf{Set}}
\newcommand{\PSh}{\mathrm{PSh}}
\newcommand{\Sh}{\mathrm{Sh}}
\newcommand{\Kar}{\operatorname{Kar}}
\newcommand{\Hom}{\operatorname{Hom}}
\newcommand{\End}{\operatorname{End}}
\newcommand{\Aut}{\operatorname{Aut}}
\newcommand{\id}{\mathrm{id}}
\newcommand{\B}{\mathbf{B}}
\newcommand{\y}{\mathbf{y}}
\newcommand{\op}{\mathrm{op}}
\newcommand{\res}[2]{#1|_{#2}}
\newcommand{\q}[1]{«\,#1\,»}
\newcommand{\lean}[1]{{\small\nolinkurl{#1}}}
\newcommand{\lp}{\href{https://github.com/LeanPool/lean-pool/tree/main/LeanPool/ToposTheoreticBridges}{\nolinkurl{ToposTheoreticBridges}}}

% Lean code blocks.
\DefineVerbatimEnvironment{leancode}{Verbatim}{fontsize=\small,frame=leftline,framerule=0.6pt,rulecolor=\color{blue!40},xleftmargin=6pt}

\theoremstyle{plain}
\newtheorem{theorem}{\thmname}[section]
\newtheorem{corollary}[theorem]{\corname}
\newtheorem{proposition}[theorem]{\propname}
\newtheorem{lemma}[theorem]{\lemname}
\newtheorem{bridge}{\bridgename}
\renewcommand{\thebridge}{\Alph{bridge}}
\theoremstyle{definition}
\newtheorem{definition}[theorem]{\defname}
\newtheorem{example}[theorem]{\exname}
\newtheorem{exercise}{\exoname}[section]
\theoremstyle{remark}
\newtheorem*{remark}{\remname}
\tcolorboxenvironment{theorem}{enhanced, breakable, colback=green!2, colframe=green!35!black, boxrule=0.6pt, arc=0pt, left=7pt, right=7pt, top=4pt, bottom=4pt, before skip=12pt, after skip=8pt}
\tcolorboxenvironment{bridge}{enhanced, breakable, colback=blue!3, colframe=blue!35!black, boxrule=0.5pt, arc=0pt, left=7pt, right=7pt, top=4pt, bottom=4pt, before skip=12pt, after skip=8pt}
\tcolorboxenvironment{definition}{enhanced, breakable, colback=black!3, colframe=black!25, boxrule=0.4pt, arc=0pt, left=7pt, right=7pt, top=4pt, bottom=4pt, before skip=10pt, after skip=6pt}
\newcommand{\entry}[1]{\par\smallskip\noindent\textbf{#1}\enspace}

\newcommand{\inwords}{\entry{En mots.}}
\newcommand{\tracename}{Trace du pipeline}
\newcommand{\gen}{\entry{Comment le générateur le trouve.}}
\newcommand{\direct}{\entry{Preuve, sans topos.}}
\newcommand{\status}{\entry{Est-ce nouveau?}}
\newcommand{\leanentry}{\entry{Lean.}}

\title{Des ponts dans les notes de Grothendieck}
\author{Michel Hua\\ \small Commutator, Paris\\ \small\texttt{michel@commutator.io} · ORCID 0009-0002-8304-0683}
\date{Brouillon, 23 novembre 2026}

\begin{document}
\maketitle
\begin{abstract}
Les \q{ponts} d'O.~Caramello lisent une même propriété d'un topos sur deux présentations de ce topos, et tirent de ces deux lectures un théorème qui relie les présentations. Un article associé \cite{bridges-generator} fait de ces ponts un générateur de théorèmes, exécuté à la main. La présente note applique ce générateur aux notes manuscrites de Grothendieck (Archives Grothendieck, Université de Montpellier). Trois arches (deux valeurs, connexe, topos de préfaisceaux d'un préordre), démontrées en Lean~4 dans le paquet \texttt{lean/bridges} de cette édition, sont lues sur trois tabliers qui reviennent dans ses notes: la catégorie simpliciale, les groupes et les préordres. Nous discutons deux précédents tirés de ces notes: d'une part un exercice du dossier 161-3, p.~26, que le générateur retrouve par l'arche de Morita et dont la forme relative aux préordres est démontrée en Lean, avec la même identification au dossier 37, p.~48; d'autre part la correspondance entre sous-topos et quotients d'une théorie, au dossier 161-2, p.~17 et 44. Un dernier tablier, tiré du dossier 161-2, le topos des pseudo-corps (anneaux réguliers au sens de von Neumann), fait l'objet d'une section: nous vérifions l'une des trois présentations qu'il en donne, nous y lisons les arches et nous répondons à une question qu'il a laissée en marge. Une autre section reprend son lemme de recollement de sous-topos (dossier 161-2, p.~48) : nous l'énonçons et le démontrons, et montrons que dans le cas où il l'applique, les anneaux de caractéristique définie, il est en défaut, de sorte qu'aucun sous-topos du topos des anneaux ne classifie ces anneaux. Une annexe recense les fichiers Lean. Une version anglaise accompagne ce texte.
\end{abstract}

\tableofcontents
\section{Introduction}\label{sec:nintro}

Un topos peut se présenter de bien des manières qui semblent sans rapport: comme les actions d'un monoïde, comme les représentations d'un carquois, comme les faisceaux sur un espace, comme les modèles d'une théorie logique. Les \emph{ponts} d'O.~Caramello \cite{caramello-unification,caramello-tst} tirent parti de cette diversité. Une propriété du topos invariante par équivalence se caractérise sur chacune de deux présentations; ces caractérisations sont les \emph{arches} du pont, et l'équivalence entre les deux présentations en est le \emph{tablier}. Mises bout à bout, les deux arches donnent un théorème qui relie les deux présentations et où le topos n'apparaît plus.

Dans un exposé aux Mines Paris \cite{caramello-mines}, Caramello a proposé de confier à une machine la combinaison systématique des arches et des tabliers. L'article associé \cite{bridges-generator} fait tourner à la main une petite version de ce générateur et la décrit comme un pipeline: un registre d'arches, démontrées en Lean~4 sous la forme \texttt{(propriété de PSh C) ↔ (condition sur C)}; un registre de tabliers, petites catégories $C$ dont les préfaisceaux forment une structure familière; enfin des étapes qui associent chaque tablier à chaque arche, calculent la condition sur $C$, la réécrivent dans la langue du domaine et la vérifient en Lean, par une preuve directe et dans la littérature. La présente note applique ce générateur aux notes manuscrites de Grothendieck conservées à l'Université de Montpellier, où l'on trouve déjà des tabliers de ce genre et des invariants lus sur deux présentations d'un même topos.

\subsection*{Plan}

La section~\ref{sec:recall} rappelle les notations et les cinq arches de l'article associé. La section~\ref{sec:morearches} donne trois arches de plus (deux valeurs, connexe, spatial), démontrées en Lean, et la section~\ref{sec:grdecks} les lit sur trois tabliers qui reviennent dans les notes de Grothendieck. La section~\ref{sec:precedents} rassemble deux précédents tirés de ces notes: l'exercice du dossier 161-3, p.~26, que le générateur retrouve et dont nous démontrons en Lean la forme relative aux préordres, avec la même identification au dossier 37, p.~48; puis les sous-topos d'une théorie, au dossier 161-2. La section~\ref{sec:pseudofields} tire du dossier 161-2 un tablier nouveau, le topos des pseudo-corps, dont Grothendieck donne trois présentations; nous en vérifions une, nous y lisons les arches et nous répondons à une question qu'il a laissée en marge. La section~\ref{sec:stacks} relie les topos classifiants aux champs de Grothendieck et propose des arches plus fines, avec un exemple tiré du dossier 146. La section~\ref{sec:nconclusion} conclut, et l'annexe~\ref{sec:nlean} donne les fichiers et les déclarations Lean.

\section{Notations et les cinq arches}\label{sec:recall}

Nous reprenons les notations de l'article associé \cite{bridges-generator}, que nous rappelons brièvement. Une petite catégorie $C$ est une \emph{forme}. Un \emph{préfaisceau} sur $C$ est un foncteur $C^{\op}\to\Set$, en Lean \texttt{Cᵒᵖ ⥤ Type u}; les préfaisceaux sur $C$ forment le topos $\PSh(C)$, que Grothendieck note $\widehat C$. Un \emph{tablier} est une forme dont les préfaisceaux constituent une structure familière. La négation $\neg A$ d'un sous-préfaisceau $A\subseteq F$ est le plus grand sous-préfaisceau de $F$ disjoint de $A$. Le topos $\PSh(C)$ est \emph{booléen} si $A\cup\neg A = F$ pour tout $A\subseteq F$ (tiers exclu), et \emph{de De Morgan} si $\neg A\cup\neg\neg A = F$ pour tout $A\subseteq F$ (tiers exclu faible). La catégorie $C$ vérifie la \emph{condition d'Ore} si toute paire de flèches de même but $f : Y\to X$, $g : Z\to X$ se complète en un carré commutatif $f\circ h = g\circ k$. L'\emph{enveloppe de Karoubi} $\Kar(C)$ ajoute un objet pour chaque flèche idempotente $e = e\circ e$, et contient $C$ comme sous-catégorie pleine. La catégorie $C$ est \emph{mince} s'il existe au plus une flèche entre deux objets quelconques.

L'article associé fait tourner son générateur sur cinq arches, démontrées en Lean~4 dans la bibliothèque \lp{} de lean-pool:
\begin{itemize}
\item \emph{booléen:} $\PSh(C)$ est booléen si et seulement si $C$ est un groupoïde;
\item \emph{De Morgan:} $\PSh(C)$ est de De Morgan si et seulement si $C$ vérifie la condition d'Ore;
\item \emph{Morita:} on a $\PSh(C)\simeq\PSh(D)$ si et seulement si $\Kar(C)$ et $\Kar(D)$ sont équivalentes;
\item \emph{points essentiels:} les foncteurs $\PSh(C)\to\Set$ qui préservent toutes les limites et toutes les colimites forment une catégorie équivalente à $\Kar(C)^{\op}$;
\item \emph{sous-topos:} les sous-topos de $\PSh(C)$ correspondent, en renversant l'ordre, aux topologies de Grothendieck sur $C$.
\end{itemize}

\section{Trois arches de plus}\label{sec:morearches}

Nous ajoutons trois arches au registre de l'article associé \cite{bridges-generator}. Nous ne les avons encore appliquées qu'aux tabliers tirés des notes de Grothendieck (\S\ref{sec:grdecks}) et à ses pseudo-corps (\S\ref{sec:pseudofields}); aucun des résultats de l'article associé n'en vient. Nous les numérotons de F à H, à la suite des cinq arches A à E de l'article associé.
\setcounter{bridge}{5}

\begin{bridge}[Deux valeurs]\label{br:twoval}
Les sous-objets de l'objet final de $\PSh(C)$ (ses \q{ouverts}) correspondent aux cribles de $C$ \cite[IV~8.4.4]{sga4}. Ainsi $\PSh(C)$ n'a pas d'autre sous-objet de l'objet final que $0$ et $1$ si et seulement si, pour deux objets quelconques $a$, $b$ de $C$, il existe une flèche $a\to b$. Pour $C$ vide, le topos est dégénéré et la condition est vérifiée par vacuité: l'énoncé porte donc sur les $C$ non vides. Lean: \lean{forall_subfunctor_terminal_eq_bot_or_eq_top_iff}.
\end{bridge}

\begin{bridge}[Connexe]\label{br:conn}
$\PSh(C)$ est connexe non vide au sens de \cite[IV~8.7~e)]{sga4} (toute partition de l'objet final a exactement une partie non vide) si et seulement si $C$ est connexe, fait rappelé comme connu dans \cite[ex.~3.6]{hora-kamio}. En termes de sous-objets de l'objet final: $C$ est non vide et les seuls sous-objets complémentés de l'objet final sont $0$ et $1$. Lean: \lean{isConnected_iff_forall_isCompl}.
\end{bridge}

\begin{bridge}[Spatial]\label{br:spatial}
$\PSh(C)$ est équivalent au topos de préfaisceaux d'un préordre, donc au topos des faisceaux sur son espace d'Alexandrov, si et seulement si $C$ est mince. Grothendieck le donne en exercice \cite[p.~26]{grothendieck-161-3}; voir le \S\ref{sec:precedents}. Lean: \lean{exists_preorder_moritaEquivalentCat_iff_isThin}, qui énonce la forme relative aux préordres; que le topos soit alors celui d'un espace relève de l'équivalence de Grothendieck, tirée du même exercice, et n'est pas formalisé ici.
\end{bridge}

Sur les deux principaux tabliers de l'article associé, les $M$-ensembles à droite et les représentations d'un carquois, elles se lisent ainsi.

\begin{center}\small
\begin{tabularx}{\linewidth}{>{\raggedright\arraybackslash}p{2.6cm}>{\raggedright\arraybackslash}p{3.2cm}>{\raggedright\arraybackslash}X>{\raggedright\arraybackslash}X}
\toprule
invariant & lecture sur $C$ & sur les $M$-ensembles à droite & sur les représentations d'un carquois\\
\midrule
deux valeurs & une flèche entre deux objets quelconques & toujours & fortement connexe\\
connexe & $C$ connexe & toujours & connexe\\
préordre (spatial) & $C$ mince & $M$ trivial & au plus un chemin entre deux sommets\\
\bottomrule
\end{tabularx}
\end{center}

\section{Les tabliers des notes de Grothendieck}\label{sec:grdecks}

Trois tabliers reviennent dans les notes de travail de Grothendieck. Voici ce qu'en disent les trois arches du \S\ref{sec:morearches} (Lean: \lean{GrothendieckDecks.lean}).
\begin{itemize}
\item La catégorie simpliciale $\Delta$, la première catégorie test de \emph{Pursuing Stacks} (dossiers 108 et 118 du fonds): les ensembles simpliciaux forment un topos à deux valeurs et connexe, puisqu'il existe une application constante entre deux objets quelconques de $\Delta$, et ce n'est pas le topos de préfaisceaux d'un préordre, car les deux sommets de $[1]$ sont deux flèches $[0]\to[1]$ (\lean{simplexCategory_twoValued}, \lean{simplexCategory_isConnected}, \lean{simplexCategory_not_spatial}).
\item Un groupe $G$ vu comme groupoïde à un objet, dont le topos de préfaisceaux est $\B G$ (dossier 161-6, p.~22, où Grothendieck ramène une question à ce cas): à deux valeurs, comme le note \cite[IV~8.4.3]{sga4} pour tout monoïde, connexe, et topos de préfaisceaux d'un préordre si et seulement si $G$ est trivial (\lean{singleObj_twoValued}, \lean{singleObj_isConnected}, \lean{singleObj_spatial_iff}).
\item Un ensemble préordonné $I$, dont le topos de préfaisceaux $\widehat I$ est $\mathrm{Top}(I)$ (dossier 37, p.~48; dossier 161-3, p.~26): à deux valeurs si et seulement si deux éléments quelconques sont comparables dans les deux sens, c'est-à-dire si et seulement si son quotient ordonné a au plus un point (\lean{preorder_twoValued_iff}).
\end{itemize}
Ce sont des corollaires immédiats; nous les consignons parce que ce sont les premières entrées du registre tirées du fonds.

\section{Précédents dans les notes de Grothendieck}\label{sec:precedents}

Nous rassemblons ici deux précédents tirés des notes de Grothendieck; les tabliers qu'on trouve dans ces notes sont déjà au registre (\S\ref{sec:grdecks}).

\subsection{Préordres et topos spatiaux: dossier 161-3, p.~26, et dossier 37, p.~48}

À la p.~26 d'une note manuscrite non datée sur les topos \cite{grothendieck-161-3}, Grothendieck énumère des conditions équivalentes portant sur un topos $E$: être équivalent à la fois au topos d'un espace et à un topos de préfaisceaux; être le topos d'un espace \q{essentiel} (un espace dont tous les points sont essentiels); être équivalent à $\widehat I$ pour un ensemble ordonné $I$; être engendré par les sous-objets de l'objet final et avoir suffisamment de points essentiels. Il demande de prouver que $\widehat C$ vérifie ces conditions si et seulement si $C$ est préordonnée, et que le topos d'un espace $X$ les vérifie si et seulement si les points essentiels de $X$ sont très denses; il note en outre des 2-équivalences entre topos spatiaux essentiels, préordres et espaces essentiels. Ce sont là des ponts au sens où nous l'entendons: le tablier est celui des préordres, et les arches sont les sous-objets de l'objet final et les points essentiels. Le générateur retrouve le premier exercice par l'arche de Morita: si $\widehat C\simeq\widehat P$ pour un préordre $P$, alors $\Kar(C)\simeq\Kar(P)\simeq P$, et $C$, sous-catégorie pleine de son enveloppe de Karoubi, est mince; réciproquement une catégorie mince est équivalente à son squelette mince, qui est un ordre (Lean: \lean{exists_preorder_moritaEquivalentCat_iff_isThin}, ci-dessous). Le second exercice, sur les espaces, n'est pas formalisé ici. La même identification figure dans ses notes pour la réédition de SGA \cite[p.~48]{grothendieck-37}: pour un ensemble ordonné $I$, les ouverts de $\widehat I$ sont les cribles de $I$, d'où $\widehat I\simeq\mathrm{Top}(I)$. Cette étape utilise que $\widehat I$ est engendré par ses sous-objets de l'objet final, ce qui est le cas parce que, pour un préordre, les préfaisceaux représentables sont des sous-objets de l'objet final.

Voici l'énoncé Lean du premier exercice, dans \lean{PreorderPresheaves.lean}. Sa preuve suit l'argument ci-dessus: \lean{karoubi_of_moritaEquivalentCat} lit l'arche de Morita, \lean{karoubiEquivPreorder} identifie un préordre à son enveloppe de Karoubi, et la réciproque passe par le squelette mince.
\begin{leancode}
variable (C : Type u) [SmallCategory C]

theorem exists_preorder_moritaEquivalentCat_iff_isThin :
    (∃ (P : Type u) (_ : Preorder P), MoritaEquivalentCat C P) ↔
      Quiver.IsThin C
\end{leancode}
L'article associé range cet énoncé dans son journal d'étalonnage, comme non publié en tant que tel \cite{bridges-generator}.

\subsection{Sous-topos et quotients d'une théorie: dossier 161-2, p.~17 et 44}

Dans ses notes d'algèbre universelle \cite[p.~17]{grothendieck-161-2}, sous le titre \q{sous-théorie d'une $\Delta$-théorie (représentable par un topos) et sous-topos}, Grothendieck prend un topos $\mathcal T$ (son $R$) et un ensemble $\Sigma$ de flèches de $\mathcal T$: les morphismes qui inversent les flèches de $\Sigma$ se factorisent par un sous-topos $\mathcal T'$ de $\mathcal T$, et tout sous-topos de $\mathcal T$ s'obtient ainsi. Si l'on prend pour $\mathcal T$ le topos classifiant d'une théorie, c'est en substance la correspondance entre quotients d'une théorie et sous-topos de son topos classifiant sur laquelle Caramello fonde ses ponts \cite{caramello-tst}. À la p.~44 du même dossier, il lit des structures comme des ouverts d'un topos classifiant.

\section{Un tablier tiré des notes de Grothendieck: les pseudo-corps}\label{sec:pseudofields}

Les précédents du \S\ref{sec:precedents} montrent Grothendieck lisant des invariants d'un topos sur deux présentations. Ses notes d'algèbre universelle contiennent aussi un tablier que le registre ne contenait pas encore: le topos des \q{pseudo-corps}, dont il donne trois présentations.

\subsection{Le tablier}

Dans le dossier 161-2 \cite{grothendieck-161-2}, pp.~38--43, Grothendieck énumère des propriétés d'anneaux et, pour chacune, un sous-topos du topos $X_{\mathrm{ann}}\simeq\widehat R$ qui classifie les anneaux, $R$ désignant la catégorie des schémas affines de type fini sur $\mathbb Z$. Sa numérotation change d'une page à l'autre: à la p.~41, l'item~12 est celui des corps, tandis qu'à la p.~42 il renumérote, de sorte que l'item~12 devient celui des anneaux compacts et l'item~13 celui des anneaux compacts réduits. Dans la suite, \q{item~12} et \q{item~13} renvoient à la numérotation de la p.~42, et nous parlons des \q{corps de la p.~41}. L'item~13 est l'\emph{anneau compact réduit}, qu'il nomme \emph{pseudo-corps}: pour tout $f$, il existe un unique idempotent $e$ tel que $ef$ soit inversible dans $eA$ et $(1-e)f=0$, de sorte que $A\simeq A_f\times A/fA$. Ce sont les anneaux commutatifs réguliers au sens de von Neumann. Il donne trois présentations du sous-topos qu'ils définissent:
\[
X_{\mathrm{ps.\,corps}}\;\simeq\;\widetilde R_{\mathrm{ps.\,corps}}\;\simeq\;\widehat{R_{\mathrm{cons}}},
\]
les faisceaux sur $R$ pour la topologie où une famille est couvrante quand l'un de ses membres a une section fibre par fibre, et les préfaisceaux sur la catégorie $R_{\mathrm{cons}}$ obtenue à partir de $R$ en inversant les flèches bijectives à extensions résiduelles triviales (pp.~41--42).

\subsection{Vérification de la présentation par préfaisceaux}

La théorie des pseudo-corps s'axiomatise par les axiomes des anneaux commutatifs et le séquent
\[
\forall a\ \exists!\,b\ \ (a^2b = a\ \wedge\ ab^2 = b):
\]
dans un anneau commutatif, le quasi-inverse $b$ de $a$ est unique, et l'idempotent de la définition de Grothendieck est $e = ab$. La théorie est donc cartésienne, et par suite de type préfaisceau: son topos classifiant est le topos des foncteurs à valeurs dans $\Set$ sur ses modèles de présentation finie. Notons $T$ l'adjoint à gauche de l'inclusion des anneaux réguliers au sens de von Neumann dans les anneaux. Cette inclusion commute aux colimites filtrantes, donc $T$ préserve la présentation finie. Réciproquement, tout pseudo-corps de présentation finie est de la forme $T(A)$, avec $A$ un anneau de présentation finie: dans une présentation finie, on remplace chaque quasi-inverse par une nouvelle variable soumise à ses deux équations. D'après Olivier \cite{olivier}, $\operatorname{Spec}T(A)$ est $\operatorname{Spec}A$ muni de sa topologie constructible, et les anneaux locaux de $T(A)$ sont les corps résiduels de $A$. D'après Pierce \cite{pierce}, un anneau régulier au sens de von Neumann est l'anneau des sections globales du faisceau de ses anneaux locaux sur son spectre, qui est un espace booléen. Ainsi, pour $u:A\to B$ de présentation finie, $T(u)$ est un isomorphisme si et seulement si $\operatorname{Spec}u$ est bijectif et toutes les extensions résiduelles sont triviales: une bijection continue entre espaces booléens est un homéomorphisme, et un morphisme de faisceaux qui est un isomorphisme sur les fibres est un isomorphisme. Les flèches que $T$ inverse sont donc exactement celles de la classe $W$ qu'inverse Grothendieck.

Reste à voir que les pseudo-corps de présentation finie s'obtiennent à partir des anneaux de présentation finie en inversant $W$; voici notre argument. La construction d'Olivier écrit $T(B)$ comme colimite filtrante de $B$-algèbres de présentation finie obtenues en adjoignant des quasi-inverses, un élément $a$ à la fois:
\[
B' = B[y]/(a^2y - a,\ ay^2 - y).
\]
Chaque morphisme $B\to B'$ de cette tour est bijectif sur les spectres, à extensions résiduelles triviales: au-dessus d'un point où $a$ est non nul dans le corps résiduel, les équations imposent $y = 1/a$, et là où $a = 0$, elles imposent $y = 0$. Ces morphismes sont donc dans $W$, de même que leurs composés. Comme $A$ est de présentation finie, un morphisme $T(A)\to T(B)$, c'est-à-dire un morphisme $A\to T(B)$, se factorise par un étage de la tour; il est donc représenté par un toit $A\to B'\leftarrow B$, avec $B\to B'$ dans $W$. La catégorie des pseudo-corps de présentation finie est ainsi la localisée de la catégorie des anneaux de présentation finie par rapport à $W$; son opposée est $R_{\mathrm{cons}}$, et le topos classifiant est $\widehat{R_{\mathrm{cons}}}$, comme l'énonce Grothendieck. Nous n'avons pas vérifié la première présentation, par la topologie $\mathrm{ps.\,corps}$, et rien dans cette section n'est formalisé.

\subsection{Les arches sur ce tablier}

\begin{itemize}
\item \emph{De Morgan:} oui, car les modèles de présentation finie d'une théorie cartésienne ont des sommes amalgamées, et donc s'amalgament (arche de De Morgan, \S\ref{sec:recall}).
\item \emph{Deux valeurs:} non, car l'énoncé $1=0$ définit un sous-objet de l'objet final qui n'est ni $0$ ni $1$.
\item \emph{Connexe:} oui, car $T(\mathbb Z)$ est un modèle de présentation finie initial, de sorte que la catégorie de ces modèles est connexe.
\item \emph{Booléen:} non, car le morphisme $T(\mathbb Z)\to T(\mathbb F_2)$ n'est pas inversible.
\item \emph{Points essentiels:} les pseudo-corps de présentation finie $T(A)$, à isomorphisme près.
\item \emph{Sous-topos:} les théories quotients de la théorie des pseudo-corps, en vertu de la dualité entre théories quotients et sous-topos \cite{caramello-tst}. Plusieurs items de la liste de Grothendieck en font partie: les corps de la p.~41, qu'il identifie aux pseudo-corps intègres et aux pseudo-corps locaux (il marque la dernière égalité \q{à prouver}, et elle est vraie, puisqu'un anneau local régulier au sens de von Neumann est un corps); les corps séparablement clos et algébriquement clos (items~14--15); les anneaux annulés par $n$ et ceux où $n$ est inversible (item~16, p.~43), qu'il marque comme sous-topos ouverts.
\end{itemize}

\subsection{Réponse à une question de la marge}

Pour l'item~12, les anneaux \emph{compacts} (tout $f$ inversible sur un idempotent $e$ et nilpotent sur $1-e$, p.~42), Grothendieck note en marge que la topologie n'est pas cohérente et demande si le topos a assez de points. À la p.~41, il avertit aussi que les sous-topos de sa liste sont déterminés par leurs points, mais que ce n'est pas clair pour leurs intersections, de sorte que vérifier une inclusion entre intersections sur les structures ensemblistes pourrait ne pas suffire. Nous levons ces deux doutes en passant par les axiomes. Chaque item de sa liste s'axiomatise, sur la signature dénombrable des anneaux, par un nombre fini de séquents géométriques. Pour les anneaux compacts, l'un de ces séquents comporte une disjonction dénombrable, puisque $(1-e)f$ est nilpotent d'un certain ordre $n$, sans borne sur $n$; c'est pourquoi la topologie n'est pas cohérente. Le topos classifiant d'une théorie géométrique dénombrable est le topos des faisceaux sur son site syntactique, catégorie dénombrable à limites finies dont la topologie est engendrée par une famille dénombrable de cribles. Or un théorème de Makkai et Reyes \cite[th.~6.2.4]{makkai-reyes} (voir aussi \cite[cor.~3.3.10]{diliberti-rogers}) affirme que le topos des faisceaux sur une catégorie dénombrable à produits fibrés, pour une topologie engendrée par une famille dénombrable de cribles, a assez de points. Le topos des anneaux compacts a donc assez de points; celui des pseudo-corps en a trivialement assez, puisque c'est un topos de préfaisceaux. Dès qu'un sous-topos $E'$ a assez de points, une inclusion $E'\subseteq E''$ se vérifie sur les points. Une intersection d'une famille dénombrable de sous-topos de sa liste classifie la théorie obtenue en réunissant leurs axiomes, qui reste dénombrable; elle a donc assez de points, et les vérifications que Grothendieck fait sur les points sont légitimes pour elle. L'argument ne couvre pas les intersections d'une famille non dénombrable de sous-topos.

\subsection{Les sous-topos ouverts, décrits par des corps}\label{sec:pseudofields-opens}

Dans la présentation par préfaisceaux, un sous-objet de l'objet final de $X_{\mathrm{ps.\,corps}}$ est une classe $U$ de pseudo-corps de présentation finie stable vers le haut par homomorphismes: si $T(A)\in U$ et s'il existe un morphisme $T(A)\to T(B)$, alors $T(B)\in U$. Le cadre des ouverts est donc celui des parties stables vers le haut du préordre \q{il existe un morphisme $T(A)\to T(B)$}, et le décrire revient à décider quand un tel morphisme existe. L'item~16 de Grothendieck (p.~43) marque deux de ces ouverts, les anneaux annulés par $n$ et ceux où $n$ est inversible. Nous décidons l'existence des morphismes par les corps résiduels, et en tirons une description de tous les ouverts. Les anneaux sont commutatifs et unitaires, \q{régulier} veut dire régulier au sens de von Neumann, et pour un anneau $A$ on pose
\[
F(A) = \{\text{corps } K \text{ tels qu'il existe un morphisme d'anneaux } A\to K\}.
\]
L'appartenance de $K$ à $F(A)$ ne dépend que des sous-corps de $K$ de type fini sur son corps premier (l'image de $A$ en engendre un quand $A$ est de type fini), de sorte que, pour $A$ de type fini, on peut voir $F(A)$ comme un ensemble de classes d'isomorphisme de corps de type fini, sans changer les inclusions entre les $F(A)$.

\paragraph{Rappels sur un anneau régulier $R$.} Nous n'utilisons que ce qui suit, élémentaire. Pour un idempotent $e$, $Re$ est un anneau d'unité $e$, et $r\mapsto er$ est un morphisme d'anneaux $R\to Re$; les parties $D(e) = \{x\in\operatorname{Spec}R : e\notin x\}$ sont ouvertes et fermées, et $\operatorname{Spec}R$ est quasi-compact. Pour un idéal premier $x$, l'anneau local $R_x$ est le corps résiduel $R/x$: si $s\notin x$ et si $s^*$ est le quasi-inverse de $s$, alors $e = ss^*$ est un idempotent tel que $s = se$, donc $e\notin x$, d'où $1-e\in x$; par suite $r/s = rs^*/1$ dans $R_x$, et $r/1 = 0$ dans $R_x$ si et seulement si $er = 0$ pour un idempotent $e\notin x$ (si $sr = 0$, prendre $e = ss^*$). Enfin, un idempotent contenu dans tous les idéaux premiers est nilpotent, donc nul.

\begin{lemma}[du local au global]\label{lem:pf-local}
Soit $A = \mathbb Z[x_1,\dots,x_n]/(f_1,\dots,f_m)$ un anneau de présentation finie et $R$ un anneau régulier (pas nécessairement de présentation finie). Si chaque corps résiduel $R/x$, $x\in\operatorname{Spec}R$, reçoit un morphisme d'anneaux depuis $A$, alors $R$ en reçoit un. Plus précisément, pour tout $x$, il existe un idempotent $e_x\notin x$ et un morphisme d'anneaux $A\to Re_x$.
\end{lemma}

\begin{proof}
Fixons $x$ et un morphisme $A\to R/x$, c'est-à-dire $a\in(R/x)^n$ tel que $f_j(a) = 0$ pour tout $j$. Comme $R\to R_x = R/x$ est surjectif, relevons $a$ en $r\in R^n$. Chaque $f_j(r)$ a une image nulle dans $R_x$, donc $e_{x,j}f_j(r) = 0$ pour un idempotent $e_{x,j}\notin x$; le produit $e_x$ des $m$ idempotents $e_{x,j}$ est un idempotent qui n'est pas dans $x$, puisque $x$ est premier. Le morphisme $\mathbb Z[x_1,\dots,x_n]\to Re_x$ qui envoie $x_i$ sur $e_xr_i$ est le composé de $x_i\mapsto r_i$ et de $R\to Re_x$; il envoie donc $f_j$ sur $e_xf_j(r) = 0$, et se factorise par $A$. C'est l'étape locale, et elle utilise les deux conditions de finitude sur $A$: un nombre fini de générateurs pour relever, un nombre fini de relations pour trouver un seul idempotent.

Les ouverts $D(e_x)$ recouvrent $\operatorname{Spec}R$, qui est quasi-compact; un nombre fini suffit, soit $D(e_1),\dots,D(e_k)$, avec des morphismes $\varphi_i\colon A\to Re_i$. Posons $\varepsilon_1 = e_1$ et $\varepsilon_i = e_i(1-e_1)\cdots(1-e_{i-1})$. Ces idempotents sont deux à deux orthogonaux, $\varepsilon_ie_i = \varepsilon_i$, et $1-\sum_i\varepsilon_i = \prod_i(1-e_i)$ est un idempotent contenu dans tout idéal premier (chaque premier évite un $e_i$, donc contient $1-e_i$): il est nul. Ainsi $R\cong\prod_iR\varepsilon_i$, et les composés $A\xrightarrow{\varphi_i}Re_i\xrightarrow{\ \cdot\varepsilon_i\ }R\varepsilon_i$ se recollent en un morphisme d'anneaux $A\to R$.
\end{proof}

On ne peut pas supprimer la finitude de $A$. Soit $P_1\sqcup P_2$ une partition des nombres premiers en deux parties infinies, et soit $A$ l'anneau engendré par $e$, des $y_q$ ($q\in P_1$) et des $z_q$ ($q\in P_2$), soumis à $e^2 = e$, $1-e = qy_q$ pour $q\in P_1$ et $e = qz_q$ pour $q\in P_2$. Les corps résiduels de $R = T(\mathbb Z)$ sont $\mathbb Q$ et les $\mathbb F_q$, et chacun reçoit un morphisme depuis $A$ ($e = 1$ en $q\in P_1$, $e = 0$ en $q\in P_2$, n'importe quoi sur $\mathbb Q$). Mais un morphisme $A\to T(\mathbb Z)$ donnerait un idempotent de $T(\mathbb Z)$ égal à $1$ exactement aux premiers de $P_1$, alors que les idempotents de $T(\mathbb Z)$ sont les parties ouvertes et fermées de $\operatorname{Spec}\mathbb Z$ muni de la topologie constructible \cite{olivier}, c'est-à-dire les ensembles finis de premiers et leurs complémentaires.

\begin{theorem}[critère]\label{th:pf-criterion}
Soient $A$ un anneau de présentation finie et $B$ un anneau quelconque. Les conditions suivantes sont équivalentes:
\begin{enumerate}
\item il existe un morphisme d'anneaux $T(A)\to T(B)$;
\item chaque corps résiduel $\kappa(\mathfrak p)$, $\mathfrak p\in\operatorname{Spec}B$, reçoit un morphisme d'anneaux depuis $A$;
\item $F(B)\subseteq F(A)$: tout corps qui reçoit un morphisme depuis $B$ en reçoit un depuis $A$.
\end{enumerate}
\end{theorem}

\begin{proof}
Par adjonction, (1) dit qu'il existe un morphisme $A\to T(B)$. D'après Olivier \cite{olivier}, $\operatorname{Spec}T(B)\to\operatorname{Spec}B$ est bijectif et le corps résiduel de $T(B)$ au point au-dessus de $\mathfrak p$ est $\kappa(\mathfrak p)$. Donc (1)$\Rightarrow$(2) en composant avec $T(B)\to\kappa(\mathfrak p)$, et (2)$\Rightarrow$(1) est le lemme~\ref{lem:pf-local} pour $R = T(B)$. Pour (2)$\Leftrightarrow$(3): les corps résiduels de $B$ sont dans $F(B)$, un morphisme $B\to K$ vers un corps se factorise par le corps résiduel en son noyau, et $F(A)$ est stable par extension de corps.
\end{proof}

La nécessité dans~(2) est immédiate; ce qu'apporte le lemme~\ref{lem:pf-local} est la suffisance, et les corps résiduels n'ont besoin d'être ni finis, ni en nombre fini. L'anneau $B$ n'a pas besoin d'être de présentation finie, et le critère dit que $T(A)\to T(B)$ existe si et seulement si $T(B)$ satisfait l'énoncé primitif positif $\sigma_A = \exists\bar x\,(f_1(\bar x) = \dots = f_m(\bar x) = 0)$, si et seulement si chaque corps résiduel de $B$ le satisfait.

\paragraph{Le treillis $D$.} Soit $D$ l'ensemble des classes $F(A)$, $A$ de présentation finie. On a $F(0) = \emptyset$, $F(\mathbb Z)$ est la classe de tous les corps, $F(A)\cap F(A') = F(A\otimes A')$ et $F(A)\cup F(A') = F(A\times A')$, la dernière égalité parce qu'un corps n'a pas d'autre idempotent que $0$ et $1$. Ainsi $D$ est un treillis distributif d'ensembles de corps. C'est le treillis de Lindenbaum des énoncés cohérents de la théorie des corps: modulo cette théorie, tout énoncé cohérent est une disjonction finie d'énoncés $\sigma_A$, et $\sigma_A\vee\sigma_{A'}$ équivaut à $\sigma_{A\times A'}$ dans un corps. Dans un pseudo-corps, $\sigma_A\vee\sigma_{A'}$ et $\sigma_{A\times A'}$ ne sont plus équivalents: $\mathbb F_2\times\mathbb F_3$ satisfait $\sigma_{\mathbb Z/6}$, mais ni $\sigma_{\mathbb Z/2}$, ni $\sigma_{\mathbb Z/3}$. C'est toute la différence entre les deux cadres ci-dessous.

\begin{theorem}[le cadre des ouverts]\label{th:pf-frame}
\begin{enumerate}
\item L'application $T(A)\mapsto F(A)$ identifie la réflexion ordonnée de la catégorie des pseudo-corps de présentation finie à l'opposé de $D$: il existe un morphisme $T(A)\to T(B)$ si et seulement si $F(B)\subseteq F(A)$.
\item Le cadre des sous-objets de l'objet final de $X_{\mathrm{ps.\,corps}}$ est isomorphe au cadre $\mathrm{Down}(D)$ de toutes les parties de $D$ stables vers le bas, ordonnées par inclusion. À une partie $S$ stable vers le bas correspond l'ouvert $U_S$, vrai dans un pseudo-corps $R$ (de présentation finie ou non) si et seulement si un membre de $S$ contient tous les corps résiduels de $R$; le sous-topos ouvert $o(U_S)$ classifie ces pseudo-corps.
\item La partie principale ${\downarrow}F(A)$ correspond à $U_A$, image du foncteur représenté par $T(A)$, qui est l'ouvert défini par $\sigma_A$. Une partie $S$ stable vers le bas est un idéal (stable par réunions finies) si et seulement si $U_S$ est déterminé par les corps résiduels, c'est-à-dire: $R\in U_S$ si et seulement si chaque corps résiduel de $R$ est dans $U_S$.
\item Le sous-topos $X_{\mathrm{corps}}$ des corps (p.~41) a pour cadre des ouverts $\mathrm{Idl}(D)$, et la restriction des ouverts de $X_{\mathrm{ps.\,corps}}$ à $X_{\mathrm{corps}}$ envoie $S$ sur l'idéal qu'elle engendre.
\end{enumerate}
\end{theorem}

\begin{proof}
(1) est le théorème~\ref{th:pf-criterion}, et chaque membre de $D$ est un $F(A)$. (2) Dans un topos de préfaisceaux $[\mathcal C,\Set]$, les sous-objets de l'objet final sont les classes d'objets de $\mathcal C$ stables vers le haut, donc, d'après~(1), les parties de $D$ stables vers le bas: on pose $S_U = \{F(B) : T(B)\in U\}$. Tout pseudo-corps $R$ est colimite filtrante de pseudo-corps de présentation finie, de sorte que $R\in U$ si et seulement si un $T(B)\in U$ s'envoie dans $R$, si et seulement si (lemme~\ref{lem:pf-local} et nécessité du théorème~\ref{th:pf-criterion}, valables pour tout $R$ régulier) tous les corps résiduels de $R$ sont dans $F(B)$ pour un $F(B)\in S_U$. (3) Le premier énoncé est~(1). Soit $S$ un idéal, et supposons chaque corps résiduel $R/x$ contenu dans un membre $F(A_x)$ de $S$. Le lemme~\ref{lem:pf-local} donne des idempotents $e_x\notin x$ et des morphismes $A_x\to Re_x$; comme dans sa preuve, un nombre fini recouvre, on les orthogonalise en $\varepsilon_1,\dots,\varepsilon_k$, et $R\cong\prod R\varepsilon_i$ reçoit un morphisme depuis $\prod A_{x_i}$, dont la classe $\bigcup F(A_{x_i})$ est dans $S$. Réciproquement, si $U_S$ est déterminé par les corps résiduels et si $F(A), F(A')\in S$, les corps résiduels de $T(A\times A')$ sont ceux de $A$ et de $A'$, tous dans $U_S$, donc $F(A)\cup F(A') = F(A\times A')\in S$. (4) $X_{\mathrm{corps}}$ classifie la théorie cohérente des corps; c'est donc le topos des faisceaux sur sa catégorie syntactique cohérente pour la topologie cohérente, et un sous-objet de l'objet final y est un crible fermé sur $1$, c'est-à-dire un ensemble d'énoncés stable par conséquence et par disjonctions finies: un idéal du treillis de Lindenbaum $D$. La restriction de $U_S$ aux corps est l'énoncé $\bigvee_{C\in S}C$, dont l'idéal est engendré par $S$.
\end{proof}

D'après~(4), les idéaux de $D$ sont les ouverts de l'espace spectral $\operatorname{Spec}D$ des filtres premiers, et chaque filtre premier est l'ensemble des membres de $D$ qui contiennent un corps donné (par le théorème de compacité: pour un filtre premier $\mathfrak p$, les énoncés $\sigma_A$, $F(A)\in\mathfrak p$, et $\neg\sigma_B$, $F(B)\notin\mathfrak p$, sont finiment satisfaisables dans les corps, puisque $F(A_1\otimes\dots\otimes A_k)\not\subseteq F(B_1\times\dots\times B_l)$). Ainsi $\operatorname{Spec}D$ est l'espace des types existentiels positifs de corps, et le cadre de $X_{\mathrm{ps.\,corps}}$ est le cadre de toutes les parties stables vers le bas de son treillis d'ouverts quasi-compacts. En termes du faisceau de Pierce \cite{pierce}: $R\in U_S$ si et seulement si les fibres du faisceau de Pierce de $R$, au-dessus de son spectre booléen, sont toutes dans un même membre de $S$. Tout ouvert non vide contient l'anneau nul, de sorte que deux ouverts non vides se rencontrent: c'est la propriété de De Morgan du topos de préfaisceaux, vue sur les ouverts.

\paragraph{Petits cas.} (a) \emph{Caractéristiques.} Envoyer $F(A)$ sur l'ensemble des caractéristiques de ses corps, c'est-à-dire sur l'image de $\operatorname{Spec}A\to\operatorname{Spec}\mathbb Z$, est un morphisme de treillis surjectif de $D$ sur le treillis des parties constructibles de $\operatorname{Spec}\mathbb Z$ (ensembles finis de premiers, et leurs complémentaires augmentés de~$0$); l'image d'un produit fibré est l'intersection des images, et la classe des corps de caractéristique dans une partie constructible $C$ est $F(\mathbb Z[1/m]\times\mathbb Z/n)$ pour des entiers sans facteur carré $m$, $n$ convenables. C'est une section, de sorte que les ouverts qui ne dépendent que de la caractéristique forment le cadre des parties stables vers le bas du treillis des constructibles de $\operatorname{Spec}\mathbb Z$. Le morphisme n'est pas injectif: pour $A = \mathbb Z[x]/((x^2-2)(x^2-3)(x^2-6))$, tout corps de caractéristique positive est dans $F(A)$, $\mathbb Q$ n'y est pas, et $\mathbb Q(\sqrt2)$ y est.

(b) \emph{L'item 16.} Dans $X_{\mathrm{ps.\,corps}}$, \q{$n = 0$} est $U_{\mathbb Z/n} = {\downarrow}V(n)$ et \q{$n$ inversible} est $U_{\mathbb Z[1/n]} = {\downarrow}D(n)$, en notant $V(n)$, $D(n)$ les classes de corps de caractéristique divisant, resp.\ ne divisant pas,~$n$. Le premier ne dépend que du radical de $n$ ($4 = 0$ et $2 = 0$ définissent le même ouvert, un pseudo-corps étant réduit), ce qui n'est pas le cas dans $\widehat R$. Leur intersection est $U_0$, le point ouvert de l'anneau nul; pour $n = 1$, c'est le N.B.\ de Grothendieck, $X_{1\text{-tors}} = X_\emptyset$, topos ponctuel. Leur réunion ${\downarrow}V(n)\cup{\downarrow}D(n)$ n'est pas le topos entier: pour $n = 2$, $\mathbb F_2\times\mathbb Q$ n'est dans aucun des deux. Mais l'idéal qu'ils engendrent est $D$ tout entier, car $V(n)\cup D(n) = F(\mathbb Z/n\times\mathbb Z[1/n])$ est l'élément maximal; en effet tout pseudo-corps se scinde en $R/nR\times R[1/n]$, le long de l'idempotent $nn^*$. Ainsi, dans $X_{\mathrm{corps}}$, les deux ouverts sont complémentaires; dans $X_{\mathrm{ps.\,corps}}$, ils sont disjoints au point près et ne recouvrent pas, et seule leur clôture par produits finis est le tout. La partie de caractéristique nulle $\bigcap_nU_{\mathbb Z[1/n]}$ a pour intérieur $U_0$ (un membre de $D$ dont tous les corps sont de caractéristique nulle a un ensemble constructible, donc vide, de caractéristiques); ainsi, comme dans $\widehat R$ (p.~47), le sous-topos des pseudo-corps sur $\mathbb Q$ n'est pas ouvert.

(c) \emph{Sur $\mathbb F_p$.} Le sous-topos ouvert $o(U_{\mathbb F_p})$ a pour cadre $\mathrm{Down}(D_p)$, $D_p = {\downarrow}F(\mathbb F_p)$, formé des classes $F(A)$ pour $A$ de présentation finie sur $\mathbb F_p$. Sur les algèbres finies étales $B = \prod_i\mathbb F_{p^{d_i}}$, qui sont leur propre $T$, il existe un morphisme $B\to B'$ si et seulement si chaque $d'_j$ est multiple d'un $d_i$; ces points forment donc le treillis des parties de $(\mathbb N_{\ge1},\mid)$ stables vers le haut et finiment engendrées, ordonnées par l'inclusion renversée. $D_p$ n'est pas déterminé par les corps algébriques: $A = \mathbb F_p[x,1/(x^p-x)]$ s'envoie dans $\mathbb F_p(t)$ et dans chaque $\mathbb F_{p^d}$, $d\ge2$, mais pas dans $\mathbb F_p$.

(d) \emph{Sur $\mathbb Q$.} Le même théorème vaut pour les $\mathbb Q$-algèbres (le lemme~\ref{lem:pf-local} est inchangé), pour le topos des pseudo-corps sur $\mathbb Q$, préfaisceaux sur les $T(A)$, $A$ de présentation finie sur $\mathbb Q$. Sur les algèbres finies étales $B = \prod K_i$, $B' = \prod K'_j$, il existe un morphisme $B\to B'$ si et seulement si chaque $K'_j$ contient un conjugué d'un $K_i$; par exemple $\mathbb Q(i)$ et $\mathbb Q(\sqrt2)$ sont incomparables, et $\mathbb Q(i)\times\mathbb Q(\sqrt2)$ s'envoie dans l'un et l'autre.

\paragraph{Ce qui est démontré, et ce qui ne l'est pas.} Le lemme~\ref{lem:pf-local} et les théorèmes~\ref{th:pf-criterion} et~\ref{th:pf-frame} sont démontrés ci-dessus à partir de la description de $\operatorname{Spec}T(B)$ par Olivier et des faits standard sur les topos de préfaisceaux et les théories cohérentes rappelés dans les preuves. Ils sont de l'édition, non de Grothendieck: les pages énoncent la présentation $\widehat{R_{\mathrm{cons}}}$ et énumèrent des ouverts, sans décrire le cadre.

\leanentry{} Le lemme~\ref{lem:pf-local} est démontré dans \href{https://github.com/Commutator-IO/GrothendieckArchivesProofs/blob/main/Grothendieck/RingTheory/Folder161_2Pseudofield.lean}{\nolinkurl{Grothendieck/RingTheory/Folder161\_2Pseudofield.lean}} (GrothendieckArchivesProofs), pour tout anneau $A$ de présentation finie et tout anneau régulier $R$, avec les corps résiduels aux idéaux maximaux, qui sont ici tous les premiers: \lean{Grothendieck.Folder161_2.exists_ringHom_of_forall_isMaximal}, l'énoncé étant \lean{OpenProblems.Findings.Folder161_2.RingHomOfResidueFields}. Le pas local est \lean{Grothendieck.Folder161_2.exists_fg_le_of_ringHom}: un morphisme $A\to R/x$ se relève en $A\to R/K$, où $K\subseteq x$ est un idéal de type fini, égal à $(g)$ pour un idempotent $g$, d'où $e_x = 1-g$. La quasi-compacité devient la finitude d'une famille des $1-g$ engendrant l'idéal unité, et l'orthogonalisation se fait un idempotent à la fois, par le théorème chinois pour les idéaux comaximaux $(g)$ et $J+(1-g)$. Le contre-exemple, le théorème~\ref{th:pf-criterion} et le théorème~\ref{th:pf-frame} ne sont pas formalisés.

\paragraph{Est-ce connu?} Pas dans ce que nous avons lu. Caramello et Johnstone \cite[p.~2 et lemmes~1.1--1.2]{caramello-johnstone} définissent, pour un anneau régulier de présentation finie $R$ (de présentation finie dans la signature munie du quasi-inverse, ce qui revient, par l'argument ci-dessus, à $R = T(A)$ avec $A$ de présentation finie), l'ensemble $\operatorname{Char}R$ des caractéristiques des corps qui reçoivent un morphisme depuis $R$, montrent qu'il est fini sans $0$ ou cofini avec $0$, et se servent des sous-objets $T_A$ de l'objet final du topos classifiant des corps définis par \q{$\operatorname{Char}\subseteq A$} pour montrer que ce topos n'est pas de De Morgan. C'est l'image de $D$ dans les constructibles de $\operatorname{Spec}\mathbb Z$ du petit cas~(a), restreinte au topos des corps. Les \emph{Universal models and definability} de Caramello \cite[\S7(ii)]{caramello-universal} reprennent ce résultat pour les caractéristiques, et ses \emph{Lattices of theories} \cite[\S7.1]{caramello-lattices} donnent le dictionnaire général qu'utilise la preuve du théorème~\ref{th:pf-frame}: les sous-objets de l'objet final d'un topos de préfaisceaux sont les sous-catégories pleines stables par cribles, et ceux d'un topos classifiant sont les énoncés géométriques modulo la théorie. Aucun des trois ne décrit le cadre de l'un ou l'autre topos. Le lemme~\ref{lem:pf-local} est l'argument de recollement des produits booléens: Burris et Sankappanavar \cite[déf.~IV.8.1 et l'exemple de Dauns et Hofmann après le cor.~IV.8.14]{burris-sankappanavar} définissent la propriété de recollement (\emph{patchwork property}) et représentent les anneaux biréguliers, donc les anneaux commutatifs réguliers, comme produits booléens d'anneaux simples, c'est-à-dire de corps, et d'un anneau nul; mais ils n'énoncent pas qu'un énoncé primitif positif est vrai dans un produit booléen si et seulement s'il l'est dans chaque facteur. Cet énoncé pourrait bien se trouver chez Burris et Werner (Trans.\ Amer.\ Math.\ Soc.\ 248, 1979) ou chez Macintyre (Fund.\ Math.\ 81, 1973), que nous n'avons pas pu lire. Nous avons interrogé TheoremSearch le 23 novembre 2026, par neuf requêtes sur les énoncés primitifs positifs dans les produits booléens et les anneaux réguliers, et sur les sous-objets de l'objet final des topos classifiants des anneaux réguliers et des corps; aucun résultat n'énonce le lemme ni la description du cadre. Le \emph{Stone Spaces} et l'article \emph{Rings, fields and spectra} de Johnstone n'ont pas été consultés.

\section{Champs et arches plus fines}\label{sec:stacks}

Les arches utilisées jusqu'ici répondent par oui ou par non: booléen ou non, connexe ou non. Les champs de Grothendieck suggèrent des arches qui calculent.

\subsection{Un topos classifiant représente un champ}

Le topos classifiant $\Set[\mathbb T]$ d'une théorie géométrique $\mathbb T$ représente, sur la 2-catégorie des topos de Grothendieck, le pseudo-foncteur qui associe à un topos $\mathcal E$ la catégorie des modèles de $\mathbb T$ dans $\mathcal E$ \cite{caramello-tst}: classifier des structures, c'est représenter un champ de modèles. Le cas le plus simple est le tablier $\B G$ d'un groupe, lu au \S\ref{sec:grdecks}: son topos classifie les $G$-torseurs, c'est donc le champ des torseurs, et les gerbes de Giraud, classées par un $H^2$ non abélien \cite{giraud}, prolongent la même idée.

\subsection{Les champs quotients, tabliers à deux présentations}

Un champ quotient $[X/G]$ a deux présentations: par le groupoïde $X\times G\rightrightarrows X$, dont le topos est celui des faisceaux $G$-équivariants sur $X$, et comme problème de modules, champ des familles d'un certain type d'objets. Lire un invariant sur les deux est un pont au sens de cette note, et les invariants qui conviennent aux champs sont numériques ou algébriques plutôt que logiques: le groupe de Picard, la cohomologie, les points. Le calcul de Mumford, $\operatorname{Pic}(\mathcal M_{1,1})\cong\mathbb Z/12$ \cite{mumford}, en est l'archétype: un invariant d'un problème de modules obtenu à partir d'une présentation explicite.

\subsection{Un exemple tiré du fonds}

Dans le dossier 146, Grothendieck dresse, pour $n$ points marqués sur la droite projective, la table des ordres de la droite tangente au dernier point, fibré en droites sur $M_{0,n}$ qui descend au quotient par $\mathfrak S_{n-1}$; il donne les ordres pour quatre points, et seulement des conjectures pour cinq et six. L'article de l'édition consacré à cette table \cite{hua-tangent} lit le groupe de Picard sur une seconde présentation du champ: $[M_{0,n}/\mathfrak S_{n-1}]$ est le champ des polynômes unitaires sans facteur carré de degré $n-1$ modulo le groupe affine, et sur cette présentation $\operatorname{Pic}[M_{0,n}/\mathfrak S_{n-1}]$ est cyclique d'ordre $(n-1)(n-2)$, engendré par la droite tangente, en caractéristique nulle et pour tout $n\ge3$. Pour $n$ pair et $n\ge6$, le groupe était connu par l'anneau de Chow entier calculé par Edidin et Hu \cite{edidin-hu}. C'est un pont avec une arche de Picard, et il donne un énoncé que nos recherches n'ont pas trouvé publié pour $n$ impair; la preuve et sa formalisation en Lean sont dans \cite{hua-tangent}.

\subsection{Les champs au sens homotopique}

\emph{À la poursuite des champs} \cite{maltsiniotis} traite d'abord des champs au sens homotopique, et ses catégories test sont des ponts d'une autre sorte encore: une catégorie $A$ est une catégorie test quand les préfaisceaux sur $A$ modélisent les types d'homotopie, et les conditions que donne Grothendieck (conditions d'asphéricité sur $A$, existence d'un intervalle séparant dans $\widehat A$) lisent une propriété homotopique de $\widehat A$ sur la catégorie $A$. Le tablier $\Delta$ du \S\ref{sec:grdecks} est la première catégorie test. Les dossiers 108 (une liste de questions), 118 (des brouillons sur la théorie de l'homotopie de $(\mathrm{Cat})$) et 134-2 (le tapuscrit du premier chapitre de \emph{Pursuing Stacks}, avec ses tables des matières) travaillent ces conditions; le \S\ref{sec:test} reprend leurs questions ouvertes.

\subsection{Un programme}

Ces exemples suggèrent pour le générateur une famille d'arches plus fines: pour un champ quotient $[X/G]$, ou plus généralement un champ muni d'une présentation explicite par un groupoïde, le groupe de Picard, la cohomologie et les points, lus sur la présentation et comparés à la description modulaire. À la différence des arches du \S\ref{sec:morearches}, elles produisent des nombres et des groupes, si bien que leurs résultats peuvent être confrontés à des tables, comme celle de Grothendieck au dossier 146, et à l'OEIS. Nous n'avons pas fait tourner le générateur avec elles; cette section ne fait que consigner le lien et un exemple traité.

\section{Anneaux de caractéristique définie : un recollement qui n'est pas un sous-topos}\label{sec:characteristic}

Aux pages~47 à~49 du dossier 161-2 \cite{grothendieck-161-2}, article~20 de sa liste, Grothendieck considère les anneaux qui sont, localement, de caractéristique $p$ pour un nombre premier $p$, ou de caractéristique $0$. Dans le topos $B = X_{\mathrm{ann}} \simeq \widehat R$ qui classifie les anneaux ($R$ les schémas affines de type fini sur $\mathbb Z$), ce sont les morphismes qui se factorisent localement par l'un de deux sous-topos :
\begin{itemize}
\item $V = X_{\mathrm{car}>0}$, réunion des sous-topos ouverts $\widehat R_{\mathbb Z/p}$, qui est ouvert : $V = o(U_V)$, où $U_V \subseteq 1$ est le crible des schémas qui s'envoient dans un $\operatorname{Spec}\mathbb F_p$ ;
\item $W = X_{\mathrm{car}\,0} \simeq \widehat{R_{\mathbb Q}}$, les $\mathbb Q$-algèbres, intersection des ouverts $o(\operatorname{Spec}\mathbb Z[1/n])$, qui n'est pas ouvert.
\end{itemize}
Il note que la condition fait intervenir une limite projective infinie, et ne donne donc pas de topos classifiant par la voie habituelle. Il cherche alors un sous-topos $Z \subseteq B$ qui classifie la condition, énonce un critère d'existence (p.~48), dont une étape est marquée \q{(à vérifier)}, vérifie le critère dans ce cas (\q{Donc OK}), et identifie le résultat au topos des préfaisceaux sur la catégorie $R_{\mathrm{car}}$ des schémas de type fini sur un corps premier, $\mathbb Q$ compris, recollés le long du schéma vide (p.~49). Nous rendons le critère précis et le démontrons, et constatons qu'il est en défaut dans ce cas : aucun sous-topos de $\widehat R$ ne classifie les anneaux de caractéristique définie. Le topos recollé existe, mais ce n'est pas un sous-topos.

\subsection{Le lemme de recollement}

Nous utilisons quatre faits classiques sur un topos $B$ \cite[\S A4.5]{elephant} :
\begin{enumerate}
\item[(F1)] les sous-topos de $B$ forment une algèbre de co-Heyting, en particulier un treillis distributif, de borne supérieure $\vee$ et de borne inférieure $\wedge$ ;
\item[(F2)] pour un sous-objet $U$ de l'objet final, un morphisme géométrique $f\colon \mathcal E\to B$ se factorise par le sous-topos ouvert $o(U)$ si et seulement si $f^*U = 1$ ; $f$ se factorise par $Y_1\wedge Y_2$ si et seulement s'il se factorise par $Y_1$ et par $Y_2$ ; enfin $o(U')\vee o(U'') = o(U'\vee U'')$ ;
\item[(F3)] se factoriser par un sous-topos est une propriété locale sur la source : si $\mathcal E$ est recouvert par des sous-objets $e_i$ de l'objet final ($\bigvee e_i = 1$) et si chaque restriction $f|_{\mathcal E/e_i}$ se factorise par $Y$, alors $f$ aussi (la condition est que $f^*$ inverse les monomorphismes denses de $Y$, et l'inversibilité est locale) ;
\item[(F4)] les sous-topos ouverts d'un sous-topos $Z\subseteq B$ sont les $Z\wedge o(U)$, $U$ sous-objet de l'objet final de $B$ (un sous-objet de l'objet final de $Z$ en est un de $B$).
\end{enumerate}
Soient $V$, $W$ deux sous-topos de $B$. Disons que $f\colon\mathcal E\to B$ est \emph{localement dans $V$ ou $W$} s'il existe des sous-objets $e'$, $e''$ de l'objet final de $\mathcal E$ avec $e'\vee e'' = 1$, $f|_{e'}$ se factorisant par $V$ et $f|_{e''}$ par $W$. Le \emph{problème de recollement} demande un sous-topos $Z$ tel que $f$ se factorise par $Z$ si et seulement si $f$ est localement dans $V$ ou $W$.

\begin{theorem}[Dossier 161-2, p.~48, rendu précis]\label{th:gluing}
Les conditions suivantes sont équivalentes :
\begin{enumerate}
\item le problème de recollement a une solution ;
\item $V\vee W$ est réunion de deux sous-topos ouverts, l'un contenu dans $V$ et l'autre dans $W$ ;
\item il existe des sous-objets $U'$, $U''$ de l'objet final de $B$ tels que
\[
W\wedge o(U')\le V,\qquad V\wedge o(U'')\le W,\qquad V\vee W\le o(U'\vee U'').
\]
\end{enumerate}
La solution est alors $Z = V\vee W$, et $Z' = V\wedge o(U')$, $Z'' = W\wedge o(U'')$ satisfont aux conditions de Grothendieck $V = Z'\vee(Z''\wedge V)$ et $W = Z''\vee(Z'\wedge W)$, avec $Z'\wedge W$ ouvert dans $W$ et $Z''\wedge V$ ouvert dans $V$.
\end{theorem}

\begin{proof}
(1)$\Rightarrow$(2). Soit $Z$ une solution. L'identité de $V$ est localement dans $V$ ou $W$, donc $V\le Z$, et de même $W\le Z$. L'inclusion de $Z$ se factorise par $Z$, elle est donc localement dans $V$ ou $W$ : $Z$ est recouvert par deux sous-topos ouverts $Z_1\le V$, $Z_2\le W$. Alors $Z = Z_1\vee Z_2\le V\vee W\le Z$. (C'est l'argument de nécessité de Grothendieck, p.~47.)

(2)$\Rightarrow$(3). D'après (F4), les deux ouverts sont $Z\wedge o(U')$ et $Z\wedge o(U'')$, avec $Z = V\vee W$. De $Z\wedge o(U')\le V$ on tire $W\wedge o(U')\le V$, et de même pour $U''$ ; et $Z = (Z\wedge o(U'))\vee(Z\wedge o(U''))\le o(U'\vee U'')$.

(3)$\Rightarrow$(1). Posons $Z = V\vee W$. Par distributivité (F1), $Z\wedge o(U') = (V\wedge o(U'))\vee(W\wedge o(U'))\le V$, et $Z\wedge o(U'')\le W$. Si $f$ est localement dans $V$ ou $W$, il se factorise par $Z$ d'après (F3). Réciproquement, si $f$ se factorise par $Z$, posons $e' = f^*U'$ et $e'' = f^*U''$ : alors $e'\vee e'' = f^*(U'\vee U'') = 1$ puisque $Z\le o(U'\vee U'')$, et $f|_{e'}$ se factorise par $Z\wedge o(U')\le V$ d'après (F2), et $f|_{e''}$ par $W$.

Pour la forme de Grothendieck : $V = V\wedge o(U'\vee U'') = (V\wedge o(U'))\vee(V\wedge o(U''))$ d'après (F1), et $V\wedge o(U'') = Z''\wedge V$ puisque $V\wedge o(U'')\le W$ ; de même pour $W$. Enfin $Z'\wedge W = W\wedge o(U')$ est ouvert dans $W$, et $Z''\wedge V$ est ouvert dans $V$.
\end{proof}

\begin{remark}
À la page~48, Grothendieck démontre la suffisance en recollant $Z'$ et $Z''$ le long de leur intersection, et marque le plongement dans $B$ qui en résulte \q{(à vérifier)}. La condition~(3) évite le recollement : la solution est la borne supérieure $V\vee W$, qui existe en tout cas. Ce que le critère ajoute, c'est la \emph{forme} des ouverts : $Z'$ doit être découpé dans $V$ par un ouvert $o(U')$ de $B$ dont la trace sur $W$ est déjà contenue dans $V$, et symétriquement. Le lemme est élémentaire une fois (F1) à (F4) admis, mais nous ne l'avons trouvé énoncé ni dans \cite{elephant} ni ailleurs.
\end{remark}

\subsection{Le cas de la caractéristique}

\begin{proposition}\label{prop:nocar}
Pour $B = \widehat R$, $V = X_{\mathrm{car}>0}$ et $W = X_{\mathrm{car}\,0}$, la condition~(3) du théorème~\ref{th:gluing} est en défaut. Aucun sous-topos de $\widehat R$ ne classifie donc les anneaux qui sont localement de caractéristique $p>0$ ou de caractéristique $0$.
\end{proposition}

\begin{proof}
Un sous-objet de l'objet final de $\widehat R$ est un crible $S$ sur $\operatorname{Spec}\mathbb Z$, c'est-à-dire un ensemble d'anneaux de présentation finie stable par les morphismes d'anneaux qui en partent, et le point de $\widehat R$ défini par un anneau $C$ est dans $o(S)$ si et seulement si un $A\in S$ admet un morphisme $A\to C$ (pour $C$ de présentation finie : si et seulement si $C\in S$). Deux observations.

(a) \emph{$W\le o(S)$ force $S$ à contenir un $\mathbb F_p$.} L'image inverse de $S$ dans $W\simeq\widehat{R_{\mathbb Q}}$ vaut $1$ si et seulement si sa valeur en $\mathbb Q$ est non vide, c'est-à-dire si $S$ contient un anneau de présentation finie $A$ muni d'un morphisme $A\to\mathbb Q$. L'image d'un tel morphisme est un sous-anneau de type fini $\mathbb Z[1/n]$ de $\mathbb Q$, donc $A$ s'envoie dans $\mathbb F_p$ pour tout premier $p\nmid n$, et $\mathbb F_p\in S$.

(b) \emph{$\mathbb F_p$ est dans $V$ et pas dans $W$, et $\mathbb Q$ est dans $W$ et pas dans $V$} : un anneau de caractéristique $p$ n'inverse aucun multiple de $p$, et $\mathbb Q$ ne s'envoie dans aucun $\mathbb F_p$.

Supposons (3). Comme $W\le o(U'\vee U'')$ et que la borne supérieure de deux cribles est leur réunion, l'argument de (a) fournit un anneau $A$ muni d'un morphisme $A\to\mathbb Q$ dans $U'$ ou dans $U''$. Si $A\in U'$, alors $W\wedge o(U') = W$ (l'image inverse de $U'$ dans $W$ vaut déjà $1$), donc $W\le V$, ce qui contredit (b) en $\mathbb Q$. Si $A\in U''$, alors $\mathbb F_p\in U''$ pour un $p$ d'après (a), donc le point $\mathbb F_p$ est dans $V\wedge o(U'')\le W$, ce qui contredit (b).
\end{proof}

Le \q{Donc OK} de Grothendieck ne tient donc pas. Ce que sa vérification établit, c'est que $V\wedge W$ est ouvert dans $V$ et dans $W$ : c'est l'ouvert $o(\emptyset)$ de l'anneau nul. Son N.B.\ de la page~48 affirme que cela suffit ; la proposition~\ref{prop:nocar} montre que non. Dans les notations du théorème~\ref{th:gluing}, le plus grand ouvert $Z'$ de $V$ de la forme voulue est $V$ lui-même, mais le plus grand ouvert $Z''$ de $W$ de cette forme est le point $o(\emptyset)\wedge W$, et $W\ne Z''\vee(Z'\wedge W)$.

\paragraph{Ce qui est connu.} La proposition~\ref{prop:nocar} est un cas immédiat d'un fait connu : être de caractéristique~$0$ n'est définissable par aucune formule géométrique dans la théorie des corps \cite[\S8.2]{caramello-flat} ; Caramello et Johnstone font la même remarque et en attribuent la première preuve à Johnstone \cite{caramello-johnstone}. En effet, les sous-objets de l'objet final d'un topos classifiant correspondent aux phrases géométriques \cite[\S7.1]{caramello-lattices}. Si un sous-topos classifiait la condition, l'ouvert $e''$ de son modèle générique serait défini par une phrase géométrique $\psi$ dans la signature des anneaux. Un corps de caractéristique~$0$ est un point de $W$ et non de $V$, donc $\psi$ y est vraie ; un corps de caractéristique~$p$ est un point de $V$ et non de $W$, donc $\psi$ y est fausse. Ainsi $\psi$ définirait la caractéristique~$0$ parmi les corps. L'étape de notre preuve selon laquelle un anneau de présentation finie qui s'envoie dans $\mathbb Q$ s'envoie dans un $\mathbb F_p$ est la forme sur les sites du lemme~1.1 de \cite{caramello-johnstone}, sur les caractéristiques des anneaux réguliers de présentation finie, et Hutzler observe qu'une borne supérieure de sous-topos ne se calcule pas localement \cite[rem.~3.5.2]{hutzler}. Ce que nous n'avons pas trouvé imprimé, c'est le critère de recollement du théorème~\ref{th:gluing}, l'énoncé en termes de la borne supérieure $V\vee W$, et le fait que le topos recollé ci-dessous ne se plonge pas dans $\widehat R$.

\paragraph{Ce qui subsiste : le topos recollé.} Un préfaisceau sur $R_{\mathrm{car}}$, somme amalgamée des $R_p$ ($p$ premier ou $0$, $R_0 = R_{\mathbb Q}$) le long de leur objet initial $\emptyset$, est une famille de préfaisceaux $F_p$ sur les $R_p$ dont les valeurs $F_p(\emptyset)$ sont identifiées. Ainsi $\widehat{R_{\mathrm{car}}}$ est le topos obtenu en recollant les $\widehat{R_p}$ le long de leur point ouvert commun $o(\emptyset)$ : c'est la \q{2-somme amalgamée sous le topos ponctuel} de l'article~21 ; en regroupant les premiers, c'est $V$ et $W$ recollés le long de $V\wedge W$, et le N.B.\ de Grothendieck identifie lui-même $V$ à $\widehat{R_{\mathrm{car}>0}}$. Les inclusions de $V$ et de $W$ se recollent en un morphisme géométrique $\widehat{R_{\mathrm{car}}}\to\widehat R$.

\begin{corollary}
Le morphisme $\widehat{R_{\mathrm{car}}}\to\widehat R$ n'est pas un plongement. Le topos $\widehat{R_{\mathrm{car}}}$ n'est donc pas équivalent, au-dessus de $\widehat R$, à un sous-topos de $\widehat R$ ; en particulier ni à $V\vee W$, ni au topos des faisceaux $\widetilde R_{\mathrm{top.car}}$ pour la topologie \q{top car} de la page~47.
\end{corollary}

\begin{proof}
S'il était un plongement d'image $Z$, les sous-topos ouverts $V$ et $W$ de $\widehat{R_{\mathrm{car}}}$, qui le recouvrent, seraient des sous-topos ouverts de $Z$ contenus dans $V$ et $W$ et le recouvrant, de sorte que la condition (2) du théorème~\ref{th:gluing} serait satisfaite, contrairement à la proposition~\ref{prop:nocar}.
\end{proof}

L'identification $X_{\mathrm{car}}\simeq\widehat{R_{\mathrm{car}}}$ de la page~49 est donc juste si $X_{\mathrm{car}}$ désigne le topos recollé, et fausse s'il désigne un sous-topos de $\widehat R$, ce qui est la façon dont la page~47 l'introduit ($X_{\mathrm{car}} = \widetilde R_{\mathrm{top.car}}$). Comme son morphisme vers $\widehat R$ n'est pas un plongement, le topos recollé ne classifie pas une propriété de l'anneau.

\paragraph{Les points.} L'article~21 affirme que le topos recollé à partir de sous-topos $Z(p)$ des $\widehat{R_p}$ a assez de points si et seulement si chaque $Z(p)$ en a assez. C'est exact : un topos recouvert par des sous-topos ouverts qui ont assez de points en a assez, et un sous-topos ouvert d'un topos qui a assez de points en a assez. Pour $\widehat{R_{\mathrm{car}}}$ lui-même, topos de préfaisceaux, c'est immédiat.

\leanentry{} Le cœur ordinal du théorème~\ref{th:gluing} est démontré dans \href{https://github.com/Commutator-IO/GrothendieckArchivesProofs/blob/main/Grothendieck/Order/Folder161_2Join.lean}{\nolinkurl{Grothendieck/Order/Folder161\_2Join.lean}} (GrothendieckArchivesProofs) : \lean{Grothendieck.Folder161_2.recollement_iff} est l'équivalence de~(2) et~(3), et \lean{Grothendieck.Folder161_2.recollement} donne la forme de Grothendieck, dans un treillis distributif quelconque et pour deux éléments quelconques tenant lieu des ouverts. L'instance sur les sous-topos n'est pas formalisée : ni Mathlib ni la bibliothèque Stacks n'ont le treillis distributif des sous-topos, le sous-topos ouvert $o(U)$ d'un sous-objet de l'objet final, ni la localité de la factorisation par un sous-topos. La proposition~\ref{prop:nocar} n'est pas formalisée.

\section{Conclusion}\label{sec:nconclusion}

Sur les tabliers des notes de Grothendieck, les trois arches du \S\ref{sec:morearches} donnent des corollaires immédiats; nous les consignons parce que ce sont les premières entrées du registre tirées du fonds. La note du dossier 161-3, p.~26, contient deux exercices qui sont des ponts sur le tablier des préordres: le générateur retrouve le premier par l'arche de Morita, et sa forme relative aux préordres est démontrée en Lean; le second, sur les espaces, n'est pas formalisé. Les pseudo-corps sont le seul tablier du registre qui relie des théories équivalentes de façon non triviale. Nous en avons vérifié à la main la présentation par préfaisceaux, non celle par la topologie $\mathrm{ps.\,corps}$, et rien de ce qui les concerne n'est formalisé; l'argument par les points couvre les intersections d'une famille dénombrable de sous-topos de sa liste, non celles d'une famille non dénombrable. La profondeur d'un théorème engendré est bornée par celle de son arche et de son tablier \cite[transparent 8]{caramello-mines}: pour lire ces notes plus profondément, il faudrait des arches sur de vrais sites, que le registre ne contient pas.

\appendix
\section{Lean}\label{sec:nlean}

Les trois arches du \S\ref{sec:morearches} sont démontrées dans le paquet Lake \texttt{lean/bridges} du dépôt de cette édition (\url{https://github.com/Commutator-IO/grothendieck-project}), qui dépend de lean-pool à une révision fixée; c'est là aussi que sont formalisés les corollaires sur les tabliers des notes de Grothendieck. Les cinq arches de \lp{}, dans lean-pool, sont démontrées en Lean~4 avec Mathlib, sans autre axiome que \texttt{propext}, \texttt{Classical.choice} et \texttt{Quot.sound}. Cette affirmation décrit le développement: les fichiers ne contiennent aucune commande \texttt{\#print axioms}, et nous ne consignons aucune vérification séparée à leur sujet. Pour \texttt{lean/bridges}, l'intégration continue du dépôt de cette édition fait un audit des axiomes à chaque modification: à la version décrite ici, il portait sur 437 déclarations, toutes limitées à \texttt{propext}, \texttt{Classical.choice} et \texttt{Quot.sound}. Les fichiers du paquet utilisés ici sont les suivants; les autres servent à l'article associé \cite{bridges-generator}.
\begin{itemize}
\item \lean{PreorderPresheaves.lean} (\lean{exists_preorder_moritaEquivalentCat_iff_isThin}, d'après l'exercice de Grothendieck, dossier 161-3, p.~26); \lean{SubterminalArches.lean} (les arches deux valeurs et connexe: \lean{forall_subfunctor_terminal_eq_bot_or_eq_top_iff}, \lean{isConnected_iff_forall_isCompl}); \lean{GrothendieckDecks.lean} (les tabliers des notes de Grothendieck: \lean{simplexCategory_twoValued}, \lean{simplexCategory_isConnected}, \lean{simplexCategory_not_spatial}, \lean{singleObj_twoValued}, \lean{singleObj_isConnected}, \lean{singleObj_spatial_iff}, \lean{preorder_twoValued_iff}).
\end{itemize}
Ne sont pas formalisés l'équivalence entre le topos de préfaisceaux d'un préordre et le topos d'un espace (pont~\ref{br:spatial}), le second exercice du dossier 161-3, p.~26 (\S\ref{sec:precedents}), et les énoncés du \S\ref{sec:pseudofields}.

\begin{center}\scriptsize\emergencystretch=3em
\begingroup\def\UrlBreaks{\do\_\do\.\do\/}
\begin{tabularx}{\linewidth}{>{\raggedright\arraybackslash}p{2.6cm}>{\raggedright\arraybackslash}X>{\raggedright\arraybackslash}p{2.9cm}>{\raggedright\arraybackslash}p{2.2cm}}
\toprule
énoncé & déclaration Lean & fichier & statut\\
\midrule
petites catégories: topos de préfaisceaux d'un préordre si et seulement si mince (Grothendieck, 161-3, p.~26) & \lean{exists_preorder_moritaEquivalentCat_iff_isThin} & \lean{PreorderPresheaves} & démontré en Lean\\
arches deux valeurs et connexe (ponts~\ref{br:twoval} et~\ref{br:conn}) & \lean{forall_subfunctor_terminal_eq_bot_or_eq_top_iff}, \lean{isConnected_iff_forall_isCompl} & \lean{SubterminalArches} & démontré en Lean\\
tabliers des notes de Grothendieck: $\Delta$, groupes, préordres (\S\ref{sec:grdecks}) & \lean{simplexCategory_twoValued}, \lean{simplexCategory_isConnected}, \lean{simplexCategory_not_spatial}, \lean{singleObj_twoValued}, \lean{singleObj_isConnected}, \lean{singleObj_spatial_iff}, \lean{preorder_twoValued_iff} & \lean{GrothendieckDecks} & démontré en Lean\\
\bottomrule
\end{tabularx}
\endgroup
\end{center}

% Bibliographie commune, en anglais : composée en anglais pour éviter les espaces françaises avant « : » et « ; ».
\begin{english}\renewcommand{\refname}{Références}
\begin{thebibliography}{99}
\bibitem{caramello-mines} O.~Caramello, \emph{Topos-empowered automatic theorem generation}, talk, CRI, Mines Paris, 17 September 2024. \url{https://www.oliviacaramello.com/Talks/MinesTalk.pdf}
\bibitem{grothendieck-161-3} A.~Grothendieck, \emph{Topos : notes manuscrites} (undated), Archives Grothendieck, Université de Montpellier, cote 161-3, \url{https://grothendieck.umontpellier.fr/}.
\bibitem{sga4} M.~Artin, A.~Grothendieck, J.-L.~Verdier, \emph{Théorie des topos et cohomologie étale des schémas} (SGA~4), tome~1, Lecture Notes in Math.~269, Springer, 1972; Exposé IV, A.~Grothendieck and J.-L.~Verdier, \emph{Topos}; re-typeset edition \url{https://www.normalesup.org/~forgogozo/SGA4/}.
\bibitem{grothendieck-161-2} A.~Grothendieck, \emph{Algèbre universelle [ou catégories] : notes manuscrites} (undated), Archives Grothendieck, Université de Montpellier, cote 161-2, \url{https://grothendieck.umontpellier.fr/}.
\bibitem{grothendieck-37} A.~Grothendieck, \emph{Rééditions SGA, EGA}, notes and typescripts (1967--1971), Archives Grothendieck, Université de Montpellier, cote 37, \url{https://grothendieck.umontpellier.fr/}.
\bibitem{elephant} P.~T. Johnstone, \emph{Sketches of an Elephant: A Topos Theory Compendium}, vols.~1--2, Oxford Logic Guides 43--44, Oxford University Press, 2002.

\bibitem{caramello-flat} O.~Caramello, Extensions of flat functors and theories of presheaf type, arXiv:1404.4610.

\bibitem{caramello-johnstone} O.~Caramello and P.~T. Johnstone, De Morgan's law and the theory of fields, arXiv:0808.1972.

\bibitem{caramello-lattices} O.~Caramello, Lattices of theories, arXiv:0905.0299.

\bibitem{hutzler} M.~Hutzler, Syntactic presentations for glued toposes and for crystalline toposes, arXiv:2206.11244.

\bibitem{olivier} J.-P.~Olivier, \emph{Anneaux absolument plats universels et épimorphismes à buts réduits}, Séminaire Samuel, Algèbre commutative 2 (1967--1968), exposé~6 (Zbl 0162.05701).
\bibitem{pierce} R.~S.~Pierce, \emph{Modules over commutative regular rings}, Mem. Amer. Math. Soc.~70, 1967 (Zbl 0152.02601).
\bibitem{makkai-reyes} M.~Makkai, G.~E.~Reyes, \emph{First Order Categorical Logic. Model-theoretical methods in the theory of topoi and related categories}, Lecture Notes in Math.~611, Springer, 1977 (Zbl 0357.18002).
\bibitem{diliberti-rogers} I.~Di~Liberti, M.~Rogers, \emph{Topoi with enough points}, arXiv:2403.15338, 2024.
\bibitem{hora-kamio} R.~Hora, Y.~Kamio, \emph{Quotient toposes of discrete dynamical systems}, J. Pure Appl. Algebra 228 (2024), no.~8, 107657; arXiv:2310.02647.
\bibitem{caramello-unification} O.~Caramello, \emph{The unification of Mathematics via Topos Theory}, arXiv:1006.3930, 2010.
\bibitem{caramello-tst} O.~Caramello, \emph{Theories, Sites, Toposes: Relating and studying mathematical theories through topos-theoretic `bridges'}, Oxford University Press, 2017.
\bibitem{giraud} J.~Giraud, \emph{Cohomologie non abélienne}, Grundlehren der math. Wissenschaften~179, Springer, 1971 (Zbl 0226.14011).
\bibitem{mumford} D.~Mumford, Picard groups of moduli problems, in \emph{Arithmetical Algebraic Geometry} (Purdue, 1963), Harper \& Row, New York, 1965, 33--81.
\bibitem{edidin-hu} D.~Edidin, Z.~Hu, The integral Chow rings of the stacks of hyperelliptic Weierstrass points, \emph{Michigan Math.\ J.} 76 (2026), no.~5, 1091--1113; arXiv:2208.00556.
\bibitem{maltsiniotis} G.~Maltsiniotis (ed.), \emph{Alexandre Grothendieck, À la poursuite des champs}, vol.~I, Documents mathématiques~20, Société mathématique de France, 2022.
\bibitem{hua-tangent} M.~Hua, \emph{The tangent line at a marked point on $M_{0,n}/\mathfrak S_{n-1}$, after a table of Grothendieck}, Hexagon preprint 2611.00189, 2026, \url{https://hexagonmath.org/2611.00189}.
\bibitem{bridges-generator} M.~Hua, \emph{Topos-theoretic bridges as a theorem generator}, 2026.
\end{thebibliography}

\end{english}
\end{document}
